
\documentclass[11pt,letterpaper]{article}
\usepackage[margin=1in]{geometry}
\sloppy\emergencystretch=4em
\usepackage{array,booktabs,tabularx,enumitem,parskip,ragged2e,titlesec,microtype,xcolor,listings,needspace,xurl,fancyhdr,graphicx,float,seqsplit}
\usepackage[T1]{fontenc}
\usepackage[utf8]{inputenc}
\usepackage{lmodern}
\usepackage[colorlinks=true,linkcolor=blue,urlcolor=blue]{hyperref}
\hypersetup{pdfauthor={Unais Ali}, pdftitle={IA 642 Midterm Memory Forensics}}
\definecolor{tblhead}{HTML}{1F4E79}
\definecolor{codebg}{HTML}{F7FAFC}
\setlength{\parindent}{0pt}
\renewcommand{\arraystretch}{1.3}
\setlist[itemize]{leftmargin=1.2em,itemsep=0.25em,topsep=0.3em}
\setlist[enumerate]{leftmargin=1.3em,itemsep=0.3em,topsep=0.3em}
\newcolumntype{Y}{>{\RaggedRight\arraybackslash}X}
\newcolumntype{P}[1]{>{\RaggedRight\arraybackslash}p{#1}}
\titleformat{name=\section,numberless}{\large\bfseries\color{tblhead}}{}{0em}{}
\titleformat{name=\subsection,numberless}{\normalsize\bfseries\color{tblhead}}{}{0em}{}
\lstset{
  basicstyle=\ttfamily\footnotesize,
  backgroundcolor=\color{codebg},
  frame=single,
  rulecolor=\color{tblhead},
  breaklines=true,
  breakatwhitespace=false,
  columns=fullflexible,
  keepspaces=true,
  showstringspaces=false,
  aboveskip=0.5em,
  belowskip=0.5em,
  xleftmargin=2pt,
  xrightmargin=2pt,
  literate=*{\\}{{\textbackslash\allowbreak}}1
}
\newcommand{\winpath}[1]{\texttt{\seqsplit{#1}}}
\pagestyle{fancy}
\fancyhf{}
\fancyfoot[L]{\footnotesize Unais Ali | IA 642 Midterm}
\fancyfoot[R]{\footnotesize Page \thepage}
\renewcommand{\headrulewidth}{0pt}
\renewcommand{\footrulewidth}{0.4pt}
\begin{document}
\begin{center}
{\LARGE\bfseries IA 642 Defensive Security}\\[0.25em]
{\Large Midterm Project: Memory Forensics}\\[0.15em]
{\large WannaCry and Stuxnet from RAM}\\[0.3em]
\end{center}

\begin{tabularx}{\textwidth}{@{}l Y l l@{}}
\toprule
\textbf{Student} & Unais Ali & \textbf{EID} & E02805019 \\
\textbf{Course} & IA 642 Defensive Security & \textbf{Week} & 05 Midterm \\
\textbf{Date} & October 5, 2026 & \textbf{Format} & Individual \\
\textbf{Filename} & \texttt{Ali\_IA642\_Midterm\_MemoryForensics.pdf} & \textbf{Tooling} & Volatility 2.6 \\
\bottomrule
\end{tabularx}

\vspace{0.5em}
\section*{Lab overview}
This midterm investigates two Windows XP classroom memory images: WannaCry (\texttt{wcry.vmem}, May 2017)
and Stuxnet (\texttt{stuxnet.vmem}, June 2011). Analysis uses Volatility~2.6 plugins equivalent to the course SIFT
\texttt{vol.py} workflow. Values below come from live plugin runs on local copies extracted from
\texttt{memory\_dump.zip}. The Symantec \textit{W32.Stuxnet Dossier} v1.3 is required reading for Parts~B and~C.
MITRE ATT\&CK software pages S0366 (WannaCry) and S0603 (Stuxnet) are used for technique mapping.

\section*{Environment and tooling}
\begin{itemize}
  \item \textbf{Evidence:} \texttt{wcry.vmem} and \texttt{stuxnet.vmem} from Canvas Week~05 Midterm /
        \texttt{memory\_dump.zip}; dossier PDF from the same module.
  \item \textbf{SIFT Workstation:} Course \texttt{SIFT-Workstation.ova} imported into VirtualBox as VM
        \texttt{SIFT-Workstation} (Ubuntu~20.04, 4~GB RAM, 4~vCPU, NAT NIC host port~2225$\rightarrow$guest~22,
        snapshot \texttt{clean-after-import}). Long-mode/64-bit enabled after OVA import.
        Volatility~2.6.1 is available inside SIFT as \texttt{/usr/local/bin/vol.py} (Module~02 workflow).
  \item \textbf{Analysis runs:} Graded plugin excerpts use Volatility~2.6 profile \texttt{WinXPSP2x86}.
        The same dumps were re-run inside SIFT with \texttt{vol.py} (pslist/pstree/psscan/cmdline/dlllist/
        filescan/modules/modscan/ssdt/driverscan/connections/sockets) and matched the host findings
        (WannaCry \texttt{tasksche.exe}/PID~1940 and \texttt{@WanaDecryptor@}/PID~740; Stuxnet triple
        \texttt{lsass.exe} PIDs). Volatility~3 was not used for graded XP net answers.
  \item \textbf{Screenshots:} Question-labeled terminal figures (A1, A2, A3, A5, B2, B5/B6, B7, B8, B10)
        plus SIFT Workstation GUI captures (\texttt{fig\_SIFT\_*}).
  \item \textbf{Safety:} Lab IOCs defanged in this report; VirusTotal used hash-only (no sample upload);
        SIFT kept on NAT / isolated for analysis.
\end{itemize}

\section*{Evidence Intake and Scope}
All timestamps in analysis are UTC unless noted as host-local from \texttt{imageinfo}.

{\footnotesize
\begin{tabularx}{\textwidth}{@{}l >{\ttfamily\scriptsize}Y r@{}}
\toprule
\textbf{Artifact} & \textbf{SHA-256} & \textbf{Size (bytes)} \\
\midrule
wcry.vmem & 76E8BE1A3761878325FDFF39A5AB1FF84922A0B18947E5268DD9175795AD2BF0 & 536870912 \\
stuxnet.vmem & 5F19FF1333FC3901FBF3FAFB50D2ADB0C495CF6D33789E5A959499A92AEEFE77 & 536870912 \\
memory\_dump.zip & 0F73E7C15F16A6073EDE0F228D33F53996953FAC63844E33AF5253478994601C & n/a \\
dossier PDF & 85F536DB77070DC2228824D27DCBDD07DB73C36E2652DCCB9F99569C3D1AEA3E & n/a \\
\bottomrule
\end{tabularx}}

\noindent\textbf{Chain of custody:} Analyst Unais Ali; received via Canvas Week~05 Midterm module;
intake hashes match final analysis copies (no transformation before Volatility~2.6 runs).
\section*{Part A: WannaCry (\texttt{wcry.vmem})}
\subsection*{A1. Identify the image and choose a profile}
\textbf{Hypothesis.} WinXPSP2x86 is the correct operational profile: \texttt{imageinfo} suggests WinXPSP2x86 and WinXPSP3x86; I instantiate WinXPSP2x86 when Image Type reports SP3. DTB/KDBG must align before process parsing.
\par\medskip\noindent\textbf{Exact command}
\begin{lstlisting}
"C:\temp\midterm\tools\vol26.exe" -f "C:\temp\midterm\dumps\wcry.vmem" imageinfo
\end{lstlisting}
\noindent\textbf{Supporting output (excerpt)}
\begin{lstlisting}
Suggested Profile(s) : WinXPSP2x86, WinXPSP3x86 (Instantiated with WinXPSP2x86)
                     AS Layer1 : IA32PagedMemory (Kernel AS)
                     AS Layer2 : FileAddressSpace (C:\temp\midterm\dumps\wcry.vmem)
                      PAE type : No PAE
                           DTB : 0x39000L
                          KDBG : 0x8054cf60L
          Number of Processors : 1
     Image Type (Service Pack) : 3
                KPCR for CPU 0 : 0xffdff000L
             KUSER_SHARED_DATA : 0xffdf0000L
           Image date and time : 2017-05-12 21:26:32 UTC+0000
     Image local date and time : 2017-05-13 02:56:32 +0530
\end{lstlisting}
\begin{figure}[H]
\centering
\includegraphics[width=0.98\textwidth]{figures/fig_A1_imageinfo.png}
\caption{A1 screenshot: imageinfo on wcry.vmem (live Volatility 2.6).}
\end{figure}
\noindent\textbf{Interpretation.} Suggested profiles include WinXPSP2x86 and WinXPSP3x86 (instantiated WinXPSP2x86). Address space is IA32PagedMemory with PAE No and Image Type (Service Pack)~3. Image date and time 2017-05-12 21:26:32~UTC marks capture near the end of encryption activity. Local time +0530 reflects host timezone configuration in the VM, not attacker geography. Verification: \texttt{pslist} succeeds under both SP2 and SP3 profile suggestions; WinXPSP2x86 retained for consistency.
\vspace{0.55em}

\subsection*{A2. The malicious lineage}
\textbf{Hypothesis.} WannaCry detonation chain runs under \texttt{explorer.exe} (interactive session), not as an SMB-spawned System child visible in this snapshot.
\par\medskip\noindent\textbf{Exact command}
\begin{lstlisting}
"C:\temp\midterm\tools\vol26.exe" -f "C:\temp\midterm\dumps\wcry.vmem" --profile WinXPSP2x86 pstree
"C:\temp\midterm\tools\vol26.exe" -f "C:\temp\midterm\dumps\wcry.vmem" --profile WinXPSP2x86 cmdline
\end{lstlisting}
\noindent\textbf{Supporting output (excerpt)}
\begin{lstlisting}
Name                                                  Pid   PPid   Thds   Hnds Time
-------------------------------------------------- ------ ------ ------ ------ ----
 0x823c8830:System                                      4      0     51    244 1970-01-01 00:00:00 UTC+0000
. 0x82169020:smss.exe                                 348      4      3     19 2017-05-12 21:21:55 UTC+0000
.. 0x8216e020:winlogon.exe                            620    348     23    536 2017-05-12 21:22:01 UTC+0000
... 0x82191658:lsass.exe                              676    620     23    353 2017-05-12 21:22:01 UTC+0000
... 0x821937f0:services.exe                           664    620     15    265 2017-05-12 21:22:01 UTC+0000
.... 0x821af7e8:svchost.exe                          1024    664     79   1366 2017-05-12 21:22:03 UTC+0000
..... 0x81f747c0:wuauclt.exe                         1768   1024      7    132 2017-05-12 21:22:52 UTC+0000
..... 0x81fea8a0:wscntfy.exe                         1168   1024      1     37 2017-05-12 21:22:56 UTC+0000
.... 0x821bea78:svchost.exe                          1152    664     10    173 2017-05-12 21:22:06 UTC+0000
.... 0x81fb95d8:svchost.exe                           260    664      5    105 2017-05-12 21:22:18 UTC+0000
.... 0x821b5230:svchost.exe                           904    664      9    227 2017-05-12 21:22:03 UTC+0000
... (22 lines total; truncated)
--- cmdline (filtered) ---
Command line : \SystemRoot\System32\smss.exe
Command line : C:\WINDOWS\system32\csrss.exe ObjectDirectory=\Windows SharedSection=1024,3072,512 Windows=On SubSystemType=Windows ServerDll=basesrv,1 ServerDll=winsrv:UserServerDllInitialization,3 ServerDll=winsrv:ConServerDllInitialization,2 ProfileControl=Off MaxRequestThreads=16
Command line : winlogon.exe
Command line : C:\WINDOWS\system32\services.exe
Command line : C:\WINDOWS\system32\lsass.exe
Command line : C:\WINDOWS\system32\svchost -k DcomLaunch
Command line : C:\WINDOWS\system32\svchost -k rpcss
Command line : C:\WINDOWS\System32\svchost.exe -k netsvcs
Command line : C:\WINDOWS\system32\svchost.exe -k NetworkService
Command line : C:\
\end{lstlisting}
\begin{figure}[H]
\centering
\includegraphics[width=0.98\textwidth]{figures/fig_A2_pstree.png}
\caption{A2 screenshot: pstree lineage on wcry.vmem.}
\end{figure}
\noindent\textbf{Interpretation.} \texttt{explorer.exe} PID~1636 lists PPID~1608 (parent not in \texttt{pslist}: exited parent, common after userinit). \texttt{tasksche.exe} PID~1940 parent~1636 at 21:22:14~UTC; \texttt{@WanaDecryptor@} PID~740 parent~1940 at 21:22:22~UTC. EPROCESS ImageFileName is 16-character truncated; full path recovered via cmdline as \winpath{C:/Intel/ivecuqmanpnirkt615/tasksche.exe} (and \texttt{@WanaDecryptor@.exe}). Root under explorer supports user-session detonation rather than worm-only remote spawn as parent.
\vspace{0.55em}

\subsection*{A3. pslist versus psscan (hidden vs exited)}
\textbf{Hypothesis.} Disagreement between \texttt{pslist} and \texttt{psscan} for WannaCry helpers reflects exited processes removed from ActiveProcessLinks, not DKOM rootkit hiding.
\par\medskip\noindent\textbf{Exact command}
\begin{lstlisting}
"C:\temp\midterm\tools\vol26.exe" -f "C:\temp\midterm\dumps\wcry.vmem" --profile WinXPSP2x86 psxview
"C:\temp\midterm\tools\vol26.exe" -f "C:\temp\midterm\dumps\wcry.vmem" --profile WinXPSP2x86 psscan
\end{lstlisting}
\noindent\textbf{Supporting output (excerpt)}
\begin{lstlisting}
Offset(P)  Name                    PID pslist psscan thrdproc pspcid csrss session deskthrd ExitTime
0x02218da0 tasksche.exe           1940 True   True   False    True   True  True    True
0x021d9da0 explorer.exe           1636 True   True   False    True   True  True    True
0x01fde308 @WanaDecryptor@         740 True   True   False    True   True  True    True
\end{lstlisting}
\begin{figure}[H]
\centering
\includegraphics[width=0.98\textwidth]{figures/fig_A3_psxview.png}
\caption{A3 screenshot: psxview showing exited helpers (pslist False / psscan True).}
\end{figure}
\noindent\textbf{Interpretation.} \texttt{psxview} shows \texttt{taskse.exe}~(536) and \texttt{taskdl.exe}~(860) False in pslist, True in psscan, with ExitTime 2017-05-12 21:26:23~UTC. Exited \texttt{@WanaDecryptor@} instances also appear only in psscan with ExitTime. Kernel mechanism: pslist walks the linked EPROCESS list; psscan searches pool tags for \_EPROCESS remnants. ExitTime populated means terminated, not stealth-hidden. On a clean system the same disagreement can occur for recently exited processes.
\vspace{0.55em}

\subsection*{A4. UTC timeline and ATT\&CK S0366 mapping}
\textbf{Hypothesis.} Encryptor and helper processes cluster just before image acquisition; map to S0366 impact and recovery-inhibition techniques.
\par\medskip\noindent\textbf{Exact command}
\begin{lstlisting}
"C:\temp\midterm\tools\vol26.exe" -f "C:\temp\midterm\dumps\wcry.vmem" --profile WinXPSP2x86 pstree
"C:\temp\midterm\tools\vol26.exe" -f "C:\temp\midterm\dumps\wcry.vmem" --profile WinXPSP2x86 psxview
\end{lstlisting}
\noindent\textbf{Supporting output (excerpt)}
\begin{lstlisting}
2017-05-12 21:22:14 UTC  tasksche.exe (1940)
2017-05-12 21:22:22 UTC  @WanaDecryptor@ (740)
2017-05-12 21:25:53 UTC  @WanaDecryptor@ (424) ExitTime (psscan)
2017-05-12 21:26:23 UTC  taskse/taskdl exited; @WanaDecryptor@ (576) ExitTime
2017-05-12 21:26:32 UTC  image acquisition (imageinfo)
\end{lstlisting}
\noindent\textbf{Interpretation.} \textbf{Memory-observed mapping (MITRE ATT\&CK S0366):} T1486 Data Encrypted for Impact (\texttt{@WanaDecryptor@}, tasksche orchestration); T1490 Inhibit System Recovery (shadow deletion helpers \texttt{taskse}/\texttt{taskdl}). Timeline shows encryption activity within about four minutes of tasksche start, immediately prior to dump.
\vspace{0.55em}

\subsection*{A5. Network artifacts and MS17-010 context}
\textbf{Hypothesis.} TCP 445 listener on System PID~4 is consistent with SMB exposure used by WannaCry propagation (CVE-2017-0144 / MS17-010); absence of connscan rows for WannaCry PIDs does not disprove spread attempts.
\par\medskip\noindent\textbf{Exact command}
\begin{lstlisting}
"C:\temp\midterm\tools\vol26.exe" -f "C:\temp\midterm\dumps\wcry.vmem" --profile WinXPSP2x86 sockets
"C:\temp\midterm\tools\vol26.exe" -f "C:\temp\midterm\dumps\wcry.vmem" --profile WinXPSP2x86 connscan
\end{lstlisting}
\noindent\textbf{Supporting output (excerpt)}
\begin{lstlisting}
Offset(V)       PID   Port  Proto Protocol        Address         Create Time
---------- -------- ------ ------ --------------- --------------- -----------
0x8223d3b8        4    138     17 UDP             192.168.56.101  2017-05-12 21:21:56 UTC+0000
0x81fb4890        4      0     47 GRE             0.0.0.0         2017-05-12 21:22:55 UTC+0000
0x81f707c0     1024   1025     17 UDP             127.0.0.1       2017-05-12 21:22:53 UTC+0000
0x81f74c18      676    500     17 UDP             0.0.0.0         2017-05-12 21:22:53 UTC+0000
0x82188658     1024    123     17 UDP             192.168.56.101  2017-05-12 21:22:53 UTC+0000
0x8219b3b8        4    445      6 TCP             0.0.0.0         2017-05-12 21:21:55 UTC+0000
0x821824b0      904    135      6 TCP             0.0.0.0         2017-05-12 21:22:03 UTC+0000
0x81ffb890        4   1027      6 TCP             0.0.0.0         2017-05-12 21:22:55 UTC+0000
0x81f6a360     1024    123     17 UDP             127.0.0.1       2017-05-12 21:22:53 UTC+0000
... (19 lines total; truncated)
--- connscan ---
Offset(P)  Local Address             Remote Address            Pid
---------- ------------------------- ------------------------- ---
\end{lstlisting}
\begin{figure}[H]
\centering
\includegraphics[width=0.98\textwidth]{figures/fig_A5_sockets.png}
\caption{A5 screenshot: sockets showing TCP/445 on System PID 4.}
\end{figure}
\noindent\textbf{Interpretation.} sockets shows TCP port 445 bound to PID~4 (System) and host IP 192.168[.]56[.]101 (VirtualBox host-only). No WannaCry PID owns connscan-visible connections at capture time. MS17-010 (EternalBlue) is the documented propagation vector for S0366; empty connection tables may reflect timing, closed sockets, or lab isolation, not evidence that the malware family lacks network capability.
\vspace{0.55em}

\subsection*{A6. Safety, defanged IOCs, and kill-switch reasoning}
\textbf{Hypothesis.} Kill-switch DNS logic (tasksche/mssecsvc) can explain global slowdown when the domain registered May~12 2017; this VM was still encrypted (unreachable DNS, race after encryptor start, or local detonation under explorer).
\par\medskip\noindent\textbf{Exact command}
\begin{lstlisting}
"C:\temp\midterm\tools\vol26.exe" -f "C:\temp\midterm\dumps\wcry.vmem" --profile WinXPSP2x86 strings / yarascan (supporting)
Manual IOC review with defanging for report.
\end{lstlisting}
\noindent\textbf{Supporting output (excerpt)}
\begin{lstlisting}
Defanged IOCs (representative):
  URL: hxxp://www[.]iuqerfsodp9ifjaposdfjhgosurijfaewrwergwea[.]com
  BTC: 12t9YDPgwueZ9NyMgw519p7AA8isjr6SMw ; 13AM4VW2dhxYgXeQepoHkHSQuy6NgaEb94
Kill-switch: malware checks domain; if resolves, exit. Registration slowed outbreak globally.
\end{lstlisting}
\noindent\textbf{Interpretation.} Host encrypted anyway, so the kill-switch likely failed from this VM (isolated lab DNS), the check raced after encryption began, or the sample path ignored the switch. Lineage under explorer supports local detonation. Limitations: strings/yarascan lack process ownership without correlation; paged-out content may be absent.
\vspace{0.55em}

\subsection*{A7. Path attribution via cmdline}
\textbf{Hypothesis.} Full filesystem paths for WannaCry binaries are recovered from EPROCESS command-line arguments, not truncated ImageFileName.
\par\medskip\noindent\textbf{Exact command}
\begin{lstlisting}
"C:\temp\midterm\tools\vol26.exe" -f "C:\temp\midterm\dumps\wcry.vmem" --profile WinXPSP2x86 cmdline
\end{lstlisting}
\noindent\textbf{Supporting output (excerpt)}
\begin{lstlisting}
tasksche.exe pid:   1940
Command line : "C:\Intel\ivecuqmanpnirkt615\tasksche.exe"

@WanaDecryptor@ pid:    740
Command line : @WanaDecryptor@.exe
\end{lstlisting}
\noindent\textbf{Interpretation.} cmdline confirms staging under \winpath{C:/Intel/.../} consistent with WannaCry dropper layout. Additional strings hits require PID correlation via cmdline, dlllist, or process-specific dumps.
\vspace{0.55em}

\subsection*{A8. ATT\&CK S0366: observed vs intel-only; attribution limits}
\textbf{Hypothesis.} Separate memory-observed techniques from intelligence-only claims; single-image analysis cannot prove nation-state attribution.
\par\medskip\noindent\textbf{Exact command}
\begin{lstlisting}
"C:\temp\midterm\tools\vol26.exe" -f "C:\temp\midterm\dumps\wcry.vmem" --profile WinXPSP2x86 pslist / cmdline / psxview (technique evidence)
\end{lstlisting}
\noindent\textbf{Supporting output (excerpt)}
\begin{lstlisting}
Observed in memory: T1486 (encryptor processes), execution chain under explorer, T1490 helpers
(taskse/taskdl exited near acquisition), SMB surface via 445 listener (environment).
Intel-only: exact exploit chain byte sequence, BTC attribution, kill-switch registration details,
campaign-wide lateral movement counts.
\end{lstlisting}
\noindent\textbf{Interpretation.} (1)~T1486 Observed (\texttt{@WanaDecryptor@}). (2)~T1490 Observed helpers (exited taskse/taskdl timing). (3)~T1059 Observed (task/execution chain). (4)~T1047 Intel-primary unless WMI artifacts are explicitly carved. Public governments attributed WannaCry to Lazarus/DPRK, but this memory image alone cannot substantiate nation-state attribution.
\vspace{0.55em}

\section*{Part B: Stuxnet (\texttt{stuxnet.vmem})}
\textit{Prerequisite reading (dossier):} Executive Summary; Stuxnet Architecture (Bypassing Behavior Blocking When Loading DLLs; Injection Technique); Installation; Load Point; Windows Rootkit Functionality.
\subsection*{B1. Identify the image and its time context}
\textbf{Hypothesis.} Same WinXPSP2x86 profile operates on the PAE-enabled Stuxnet image; compare AS layer and capture timestamps to wcry.
\par\medskip\noindent\textbf{Exact command}
\begin{lstlisting}
"C:\temp\midterm\tools\vol26.exe" -f "C:\temp\midterm\dumps\stuxnet.vmem" imageinfo
\end{lstlisting}
\noindent\textbf{Supporting output (excerpt)}
\begin{lstlisting}
Suggested Profile(s) : WinXPSP2x86, WinXPSP3x86 (Instantiated with WinXPSP2x86)
                     AS Layer1 : IA32PagedMemoryPae (Kernel AS)
                     AS Layer2 : FileAddressSpace (C:\temp\midterm\dumps\stuxnet.vmem)
                      PAE type : PAE
                           DTB : 0x319000L
                          KDBG : 0x80545ae0L
          Number of Processors : 1
     Image Type (Service Pack) : 3
                KPCR for CPU 0 : 0xffdff000L
             KUSER_SHARED_DATA : 0xffdf0000L
           Image date and time : 2011-06-03 04:31:36 UTC+0000
     Image local date and time : 2011-06-03 00:31:36 -0400
\end{lstlisting}
\noindent\textbf{Interpretation.} Stuxnet: WinXPSP2x86 instantiated; AS Layer IA32PagedMemoryPae (PAE Yes) versus wcry No~PAE. Image 2011-06-03 04:31:36~UTC; local 2011-06-03 00:31:36 $-$0400. Process \texttt{pslist} shows boot-era timestamps (2010-10-29) on some processes versus infection-day 2011-06-03 on counterfeit lsass, consistent with a snapshot/resumed VM rather than a fresh boot on infection day alone.
\vspace{0.55em}

\subsection*{B2. Identifying genuine vs counterfeit lsass.exe}
\textbf{Hypothesis.} Exactly one lsass (PID~680) matches Windows XP expectations: parent winlogon, full LSA modules, clean malfind, and socket ownership.
\par\medskip\noindent\textbf{Exact command}
\begin{lstlisting}
"C:\temp\midterm\tools\vol26.exe" -f "C:\temp\midterm\dumps\stuxnet.vmem" --profile WinXPSP2x86 pslist
"C:\temp\midterm\tools\vol26.exe" -f "C:\temp\midterm\dumps\stuxnet.vmem" --profile WinXPSP2x86 pstree
\end{lstlisting}
\noindent\textbf{Supporting output (excerpt)}
\begin{lstlisting}
0x81e70020 lsass.exe  680  624  ... 2010-10-29 17:08:54 UTC
0x81c498c8 lsass.exe  868  668  ... 2011-06-03 04:26:55 UTC
0x81c47c00 lsass.exe 1928  668  ... 2011-06-03 04:26:55 UTC
\end{lstlisting}
\begin{figure}[H]
\centering
\includegraphics[width=0.98\textwidth]{figures/fig_B2_pslist.png}
\caption{B2 screenshot: pslist showing three lsass.exe instances.}
\end{figure}
\noindent\textbf{Interpretation.} Evidence bundle beyond slide~22: (1)~Parent PID~624 winlogon for 680 vs 668 services.exe for counterfeits. (2)~dlllist counts 57~/~8~/~28 (B5). (3)~malfind: 0 regions on 680 vs injected MZ on 868/1928 (B7). (4)~UDP 500/4500 sockets only on PID~680 (B4). XP norm: one legitimate lsass under winlogon.
\vspace{0.55em}

\subsection*{B3. Infection-day process context}
\textbf{Hypothesis.} Distinguish Stuxnet artifacts from analysis tooling captured in the same memory image.
\par\medskip\noindent\textbf{Exact command}
\begin{lstlisting}
"C:\temp\midterm\tools\vol26.exe" -f "C:\temp\midterm\dumps\stuxnet.vmem" --profile WinXPSP2x86 pslist
"C:\temp\midterm\tools\vol26.exe" -f "C:\temp\midterm\dumps\stuxnet.vmem" --profile WinXPSP2x86 cmdline
\end{lstlisting}
\noindent\textbf{Supporting output (excerpt)}
\begin{lstlisting}
Offset(V)  Name                    PID   PPID   Thds     Hnds   Sess  Wow64 Start                          Exit
---------- -------------------- ------ ------ ------ -------- ------ ------ ------------------------------ ------------------------------
0x823c8830 System                    4      0     59      403 ------      0
0x820df020 smss.exe                376      4      3       19 ------      0 2010-10-29 17:08:53 UTC+0000
0x821a2da0 csrss.exe               600    376     11      395      0      0 2010-10-29 17:08:54 UTC+0000
0x81da5650 winlogon.exe            624    376     19      570      0      0 2010-10-29 17:08:54 UTC+0000
0x82073020 services.exe            668    624     21      431      0      0 2010-10-29 17:08:54 UTC+0000
0x81e70020 lsass.exe               680    624     19      342      0      0 2010-10-29 17:08:54 UTC+0000
0x823315d8 vmacthlp.exe            844    668      1       25      0      0 2010-10-29 17:08:55 UTC+0000
0x81db8da0 svchost.exe             856    668     17      193      0      0 2010-10-29 17:08:55 UTC+0000
0x81e61da0 svchost.exe             940    668     13      312      0      0 2010-10-29 17:08:55 UTC+0000
0x822843e8 svchost.exe            1032    668     61     1169      0      0 2010-10-29 17:08:55 UTC+0000
0x81e18b28 svchost.exe            1080    668      5       80      0      0 2010-10-29 17:08:55 UTC+0000
... (33 lines total; truncated)
\end{lstlisting}
\noindent\textbf{Interpretation.} Infection-day (2011-06-03): counterfeit lsass 868/1928, wmiprvse activity, short-lived cmd/ipconfig. Procmon PID~660 indicates Process Monitor was running (lab capture with live monitoring); treat as environmental artifact, not Stuxnet TTP.
\vspace{0.55em}

\subsection*{B4. Socket ownership corroboration}
\textbf{Hypothesis.} Network bindings on lsass differentiate the genuine instance; supportive evidence, not sole proof.
\par\medskip\noindent\textbf{Exact command}
\begin{lstlisting}
"C:\temp\midterm\tools\vol26.exe" -f "C:\temp\midterm\dumps\stuxnet.vmem" --profile WinXPSP2x86 sockets
\end{lstlisting}
\noindent\textbf{Supporting output (excerpt)}
\begin{lstlisting}
Offset(V)       PID   Port  Proto Protocol        Address         Create Time
---------- -------- ------ ------ --------------- --------------- -----------
0x81dc2008      680    500     17 UDP             0.0.0.0         2010-10-29 17:09:05 UTC+0000
0x82061c08        4    445      6 TCP             0.0.0.0         2010-10-29 17:08:53 UTC+0000
0x82294aa8      940    135      6 TCP             0.0.0.0         2010-10-29 17:08:55 UTC+0000
0x821a5008      188   1025      6 TCP             127.0.0.1       2010-10-29 17:09:09 UTC+0000
0x81cb3d70     1080   1141     17 UDP             0.0.0.0         2010-10-31 16:36:16 UTC+0000
0x81da4d18      680      0    255 Reserved        0.0.0.0         2010-10-29 17:09:05 UTC+0000
0x81fdbe98     1032    123     17 UDP             127.0.0.1       2011-06-03 04:25:47 UTC+0000
... (14 lines total; truncated)
\end{lstlisting}
\noindent\textbf{Interpretation.} Only PID~680 owns UDP~500 and UDP~4500 (IPsec/IKE-related legitimate lsass behavior). Counterfeit PIDs 868 and 1928 have no parallel socket rows. Corroborates B2; a socket-only heuristic would fail if malware hooked networking elsewhere.
\vspace{0.55em}

\subsection*{B5. DLL inventory counts and DFIR-07 critique}
\textbf{Hypothesis.} Count dlllist module rows (0x-prefixed bases), not \texttt{wc -l} on raw output (counts header lines).
\par\medskip\noindent\textbf{Exact command}
\begin{lstlisting}
"C:\temp\midterm\tools\vol26.exe" -f "C:\temp\midterm\dumps\stuxnet.vmem" --profile WinXPSP2x86 dlllist -p 680
"C:\temp\midterm\tools\vol26.exe" -f "C:\temp\midterm\dumps\stuxnet.vmem" --profile WinXPSP2x86 dlllist -p 868
"C:\temp\midterm\tools\vol26.exe" -f "C:\temp\midterm\dumps\stuxnet.vmem" --profile WinXPSP2x86 dlllist -p 1928
\end{lstlisting}
\noindent\textbf{Supporting output (excerpt)}
\begin{lstlisting}
dlllist 0x-base row counts: PID 680 -> 57 ; PID 868 -> 8 ; PID 1928 -> 28.
(Critique: dlllist | wc -l includes banners/separators; invalid for module counts.)
\end{lstlisting}
\noindent\textbf{Interpretation.} Counterfeits load far fewer modules because hollowed/injected images do not complete full LSA initialization (dossier Injection Technique). Genuine 680 loads LSASRV, SAMSRV, msv1\_0, kerberos, and related packages (B6).
\vspace{0.55em}

\subsection*{B6. LSA security modules on genuine lsass only}
\textbf{Hypothesis.} LSA authentication stack DLLs should appear only in the real lsass process.
\par\medskip\noindent\textbf{Exact command}
\begin{lstlisting}
"C:\temp\midterm\tools\vol26.exe" -f "C:\temp\midterm\dumps\stuxnet.vmem" --profile WinXPSP2x86 dlllist -p 680
\end{lstlisting}
\noindent\textbf{Supporting output (excerpt)}
\begin{lstlisting}
************************************************************************
lsass.exe pid:    680
Command line : C:\WINDOWS\system32\lsass.exe
Service Pack 3

Base             Size  LoadCount Path
---------- ---------- ---------- ----
0x01000000     0x6000     0xffff C:\WINDOWS\system32\lsass.exe
0x7c900000    0xaf000     0xffff C:\WINDOWS\system32\ntdll.dll
0x7c800000    0xf6000     0xffff C:\WINDOWS\system32\kernel32.dll
0x77dd0000    0x9b000     0xffff C:\WINDOWS\system32\ADVAPI32.dll
0x77e70000    0x92000     0xffff C:\WINDOWS\system32\RPCRT4.dll
0x77fe0000    0x11000     0xffff C:\WINDOWS\system32\Secur32.dll
0x75730000    0xb5000     0xffff C:\WINDOWS\system32\LSASRV.dll
0x71b20000    0x12000     0xffff C:\WINDOWS\system32\MPR.dll
... (64 lines total; truncated)
\end{lstlisting}
\begin{figure}[H]
\centering
\includegraphics[width=0.98\textwidth]{figures/fig_B5_dlllist_680.png}
\caption{B5/B6 screenshot: dlllist on genuine lsass PID 680.}
\end{figure}
\noindent\textbf{Interpretation.} PID~680 maps LSASRV.dll, SAMSRV.dll, cryptdll.dll, and related security packages. Counterfeit instances lack this full stack, consistent with process injection posing as lsass.exe.
\vspace{0.55em}

\subsection*{B7. malfind: injected regions and disassembly caveat}
\textbf{Hypothesis.} PAGE\_EXECUTE\_READWRITE private VADs with MZ headers indicate injection; disassembly of PE header bytes is not shellcode.
\par\medskip\noindent\textbf{Exact command}
\begin{lstlisting}
"C:\temp\midterm\tools\vol26.exe" -f "C:\temp\midterm\dumps\stuxnet.vmem" --profile WinXPSP2x86 malfind -p 680
"C:\temp\midterm\tools\vol26.exe" -f "C:\temp\midterm\dumps\stuxnet.vmem" --profile WinXPSP2x86 malfind -p 868
"C:\temp\midterm\tools\vol26.exe" -f "C:\temp\midterm\dumps\stuxnet.vmem" --profile WinXPSP2x86 malfind -p 1928
\end{lstlisting}
\noindent\textbf{Supporting output (excerpt)}
\begin{lstlisting}
Process: lsass.exe Pid: 868 Address: 0x80000
Vad Tag: Vad  Protection: PAGE_EXECUTE_READWRITE
Flags: Protection: 6

0x00080000  4d 5a 90 00 03 00 00 00 04 00 00 00 ff ff 00 00   MZ..............
0x00080010  b8 00 00 00 00 00 00 00 40 00 00 00 00 00 00 00   ........@.......
0x00080020  00 00 00 00 00 00 00 00 00 00 00 00 00 00 00 00   ................
0x00080030  00 00 00 00 00 00 00 00 00 00 00 00 08 01 00 00   ................

0x00080000 4d               DEC EBP
0x00080001 5a               POP EDX
0x00080002 90               NOP
0x00080003 0003             ADD [EBX], AL
0x00080005 0000             ADD [EAX], AL
0x00080007 000400           ADD [EAX+EAX], AL
... (82 lines total; truncated)
--- PID 1928 ---
Process: lsass.exe Pid: 1928 Address: 0x80000
Vad Tag: Vad  Protection: PAGE_EXECUTE_READWRITE
Flags: Protection: 6

0x00080000  4d 5a 90 00 03 00 00 00 04 00 00 00 ff ff 00 00   MZ..............
0x00080010  b8 00 00 00 00 00 00 00 40 00 00 00 00 00 00 00   ........@.......
0x00080020  00 00 00 00 00 00 00 00 00 00 00 00 00 00 00 00   ................
0x00080030  00 00 00 00 00 00 00 00 00 00 00 00 08 01 00 00   ................

0x00080000 4d               DEC EBP
0x00080001 5a               POP EDX
0x00080002 90               NOP
0x00080003 0003             ADD [EBX], AL
... (197 lines total; truncated)
\end{lstlisting}
\begin{figure}[H]
\centering
\includegraphics[width=0.98\textwidth]{figures/fig_B7_malfind_868.png}
\caption{B7 screenshot: malfind RWX/MZ region in counterfeit lsass PID 868.}
\end{figure}
\noindent\textbf{Interpretation.} PID~680: 0 malfind regions (control). PID~868: 2 regions (0x80000, 0x1000000). PID~1928: 5 regions. Signals: PAGE\_EXECUTE\_READWRITE, no file backing, MZ \texttt{4d 5a}. Volatility disassembly showing DEC EBP; POP EDX; NOP is PE header misinterpreted, not Stuxnet shellcode. Control comparison on genuine lsass is required.
\vspace{0.55em}

\subsection*{B8. Process dump naming and hash correlation}
\textbf{Hypothesis.} \texttt{malfind -D} dumps use \texttt{process.<EPROCESS>.<VAD start>.dmp}; six unique SHA-256 values across regions.
\par\medskip\noindent\textbf{Exact command}
\begin{lstlisting}
"C:\temp\midterm\tools\vol26.exe" -f "C:\temp\midterm\dumps\stuxnet.vmem" --profile WinXPSP2x86 malfind -D dumps/ -p 868
"C:\temp\midterm\tools\vol26.exe" -f "C:\temp\midterm\dumps\stuxnet.vmem" --profile WinXPSP2x86 malfind -D dumps/ -p 1928
sha256sum process.*.dmp
\end{lstlisting}
\noindent\textbf{Supporting output (excerpt)}
\begin{lstlisting}
e97d61f7393ac5838a1800f3e9aa22c6205f4d7e2bde494573d35c57bc9b7819  process.0x81c47c00.0x1000000.dmp
163b7da37df4ae6dafbfb5bf88b319dabf7846cee73d4192c6a7593e835857a8  process.0x81c47c00.0x680000.dmp
abce3e79e26b5116fe7f3d40d21eaa4c8563e433b6086f9ec07c2925593f69dc  process.0x81c47c00.0x6f0000.dmp
2b2945f7cc7cf5b30ccdf37e2adbb236594208e409133bcd56f57f7c009ffe6d  process.0x81c47c00.0x80000.dmp
10f07b9fbbc6a8c6dc4abf7a3d31a01e478accd115b33ec94fe885cb296a3586  process.0x81c47c00.0x870000.dmp
a4b4b29f0df45283b629203b080c09ddb5bc6eb4cd8e9b725f75121a8b7e728e  process.0x81c498c8.0x1000000.dmp
2b2945f7cc7cf5b30ccdf37e2adbb236594208e409133bcd56f57f7c009ffe6d  process.0x81c498c8.0x80000.dmp
\end{lstlisting}
\begin{figure}[H]
\centering
\includegraphics[width=0.98\textwidth]{figures/fig_B8_hashes.png}
\caption{B8 screenshot: SHA-256 hashes of malfind -D region dumps.}
\end{figure}
\noindent\textbf{Interpretation.} Seven dump files, six unique hashes (0x80000 shared across both PIDs). 0x1000000 aligns with typical ImageBase / injection layout described in the dossier Injection Technique section.
\vspace{0.55em}

\subsection*{B9. VirusTotal hash lookup and OPSEC}
\textbf{Hypothesis.} Hash-not-upload preserves Meridian/lab OPSEC; VT labels vary (Stuxnet vs Duqu families).
\par\medskip\noindent\textbf{Exact command}
\begin{lstlisting}
Browser: VirusTotal GUI hash search (no upload) on B8 SHA-256 values.
\end{lstlisting}
\noindent\textbf{Supporting output (excerpt)}
\begin{lstlisting}
2b2945f7cc7cf5b30ccdf37e2adbb236594208e409133bcd56f57f7c009ffe6d -> 61/71 worm.stuxnet
10f07b9fbbc6a8c6dc4abf7a3d31a01e478accd115b33ec94fe885cb296a3586 -> 48/70 trojan.stuxnet / duqu
\end{lstlisting}
\noindent\textbf{Interpretation.} Divergent vendor labels reflect shared code artifacts between Stuxnet and Duqu clusters. Reporting uses hash lookup only (no sample upload) for OPSEC.
\vspace{0.55em}

\subsection*{B10. ldrmodules, drivers, and registry timeline}
\textbf{Hypothesis.} Stuxnet load-order anomalies and kernel drivers precede counterfeit lsass creation.
\par\medskip\noindent\textbf{Exact command}
\begin{lstlisting}
"C:\temp\midterm\tools\vol26.exe" -f "C:\temp\midterm\dumps\stuxnet.vmem" --profile WinXPSP2x86 ldrmodules -p 1928
"C:\temp\midterm\tools\vol26.exe" -f "C:\temp\midterm\dumps\stuxnet.vmem" --profile WinXPSP2x86 driverscan
"C:\temp\midterm\tools\vol26.exe" -f "C:\temp\midterm\dumps\stuxnet.vmem" --profile WinXPSP2x86 svcscan
\end{lstlisting}
\noindent\textbf{Supporting output (excerpt)}
\begin{lstlisting}
Pid      Process              Base       InLoad InInit InMem MappedPath
-------- -------------------- ---------- ------ ------ ----- ----------
    1928 lsass.exe            0x00080000 False  False  False
    1928 lsass.exe            0x7c900000 True   True   True  \WINDOWS\system32\ntdll.dll
    1928 lsass.exe            0x773d0000 True   True   True  \WINDOWS\WinSxS\x86_Microsoft.Windows.Common-Controls_6595b64144ccf1df_6.0.2600.5512_x-ww_35d4ce83\comctl32.dll
    1928 lsass.exe            0x77f60000 True   True   True  \WINDOWS\system32\shlwapi.dll
    1928 lsass.exe            0x771b0000 True   True   True  \WINDOWS\system32\wininet.dll
    1928 lsass.exe            0x77a80000 True   True   True  \WINDOWS\system32\crypt32.dll
    1928 lsass.exe            0x77fe0000 True   True   True  \WINDOWS\system32\secur32.dll
    1928 lsass.exe            0x77c00000 True   True   True  \WINDOWS\system32\version.dll
    1928 lsass.exe            0x01000000 True   False  True
... (31 lines total; truncated)
--- driverscan ---
MRxCls / MRxNet present
Registry Services MRxCls/MRxNet LastWrite 2011-06-03 04:26:47 UTC
\end{lstlisting}
\begin{figure}[H]
\centering
\includegraphics[width=0.98\textwidth]{figures/fig_B10_ldrmodules.png}
\caption{B10 screenshot: ldrmodules on counterfeit lsass PID 1928.}
\end{figure}
\noindent\textbf{Interpretation.} ldrmodules PID~1928: 0x80000 InLoad/InInit/InMem False (unlinked injection). Impossible KERNEL32.DLL.ASLR-style name matches dossier section ``Bypassing Behavior Blocking When Loading DLLs.'' Registry Services MRxCls/MRxNet last updated 2011-06-03 04:26:47~UTC, before counterfeit lsass at 04:26:55~UTC (driver-first installation sequence; dossier Load Point / Windows Rootkit Functionality).
\vspace{0.55em}


\section*{Part C: The Dossier on Trial}
\subsection*{C1. Claims against evidence}
At least six testable claims from four or more dossier sections. Verdicts use assignment vocabulary:
\textit{Confirmed in memory} / \textit{Consistent but not proven} / \textit{Not observable in RAM} / \textit{Contradicted}.

{\footnotesize
\begin{tabularx}{\textwidth}{@{}c >{\RaggedRight\arraybackslash}p{1.7in} >{\RaggedRight\arraybackslash}p{1.15in} c >{\RaggedRight\arraybackslash}Y@{}}
\toprule
\# & Dossier section / quoted claim & Memory verdict & Conf. & Basis \\
\midrule
1 & \textbf{Injection Technique:} ``The potential target processes for the injection are as follows: Lsass.exe\ldots'' & Confirmed in memory & High & malfind MZ/RWX on PIDs 868/1928; ldrmodules unlinked 0x80000 \\
2 & \textbf{Installation / Export 16:} ``injects itself into the services.exe process\ldots'' & Consistent but not proven & Med & PPID 668 services.exe + malfind MZ at 0x13f0000 in services.exe (C2) \\
3 & \textbf{Load Point:} ``Mrxcls.sys is a driver that allows Stuxnet to be executed every time an infected system boots\ldots'' & Confirmed in memory & High & driverscan/modules list mrxcls.sys; Services MRxCls key \\
4 & \textbf{Windows Rootkit Functionality / Table~4:} Resource 242 ``MRxnet.sys rootkit driver'' & Confirmed in memory & High & driverscan/modules list mrxnet.sys / MRxNet \\
5 & \textbf{Bypassing Behavior Blocking:} crafted names ``KERNEL32.DLL.ASLR.[HEXADECIMAL]'' & Confirmed in memory & High & ldrmodules shows ASLR-style impossible module strings \\
6 & \textbf{Executive Summary / WinCC:} copies/executes on systems running a WinCC database server; PLC rootkit claims & Not observable in RAM & N/A & No WinCC/PLC artifacts in this guest VMEM \\
7 & \textbf{Injection Technique:} suspended-process / template PE injection into trusted processes & Confirmed in memory & High & Multiple lsass ImageFileName with divergent dlllist/malfind (680 vs 868/1928) \\
\bottomrule
\end{tabularx}}

\noindent\textbf{Memory more precise than dossier:} exact PIDs 680/868/1928, dll counts 57/8/28, and second-resolution
driver LastWrite vs counterfeit create times. \textbf{Claim~6} is the required not-observable case (ICS/PLC/WinCC
scope needs disk, network, or PLC evidence outside this VMEM).

\subsection*{C2. Hypothesis: services.exe as injector}
\textbf{Hypothesis (stated before testing):} Counterfeit lsass processes (PPID~668 \texttt{services.exe}) imply
\texttt{services.exe} participated in injection or staging, matching Installation (``injects itself into the services.exe process'')
and Load Point (services.exe as injection target).

\textbf{Predictions if true:} (1)~\texttt{malfind -p 668} shows private RWX/MZ regions;
(2)~if false (parentage only), malfind on 668 is empty while counterfeit lsass still show injection.

\noindent\textbf{Exact command}
\begin{lstlisting}
"C:\temp\midterm\tools\vol26.exe" -f "C:\temp\midterm\dumps\stuxnet.vmem" --profile WinXPSP2x86 malfind -p 668
"C:\temp\midterm\tools\vol26.exe" -f "C:\temp\midterm\dumps\stuxnet.vmem" --profile WinXPSP2x86 pslist | findstr services
\end{lstlisting}
\noindent\textbf{Observed / verdict:} malfind on services.exe reveals MZ header at 0x13f0000 (RWX private VAD).
Verdict: \textit{Consistent but not proven} that services.exe performed the hollowing (could be residual payload
or related tooling). Parentage alone remains circumstantial; positive malfind elevates services.exe as a
loader/orchestrator candidate. Drivers already registered (B10), so user-mode hollowing may occur without a
user-visible reboot in a resumed VM.

\subsection*{C3. Source critique}
The Symantec \textit{W32.Stuxnet Dossier} v1.3 (November 2010) predates this image's capture date (2011-06-03)
and is primary-sourced technical intelligence from static reverse engineering and field telemetry.
\textbf{Strengths:} precise injection narrative, driver names (MRxCls/MRxNet), trusted-process targets, WinCC/PLC context.
\textbf{Limits:} publication lag relative to later variants; vendor-centric attribution language; air-gapped PLC spread
cannot be confirmed or denied from one XP VMEM alone.
\textbf{Untestable from RAM alone (example):} PLC project modification / ``PLC rootkit'' claims require PLC/project
backups or engineering-station disk artifacts, not guest RAM alone.
Analyst treats the dossier as a hypothesis generator; Volatility outputs are falsifiable tests
(\texttt{malfind}, \texttt{ldrmodules}, \texttt{dlllist}). Conflicts require tiering: memory-observed $>$ correlated disk/network $>$ dossier-only.


\section*{Part D: From Evidence to Action}
\subsection*{D1. Findings memo to Meridian's CISO}
\begin{quote}
\small
\textbf{TO:} Meridian Energy CISO and SOC Lead\\
\textbf{FROM:} Unais Ali, Memory Forensics Analyst (IA~642)\\
\textbf{SUBJECT:} Dual-incident memory triage: WannaCry lab image and Stuxnet reference image\\
\textbf{DATE:} October 5, 2026
\end{quote}

\subsubsection*{Executive summary (non-technical)}
We analyzed two public reference memory images to prove Meridian can find malware in RAM when disk
artifacts are missing, encrypted, or hidden. The WannaCry image shows ransomware already encrypting a
Windows XP lab machine in May~2017, with encryption helpers finishing seconds before the memory capture.
Disk-only triage would have seen encrypted files but would likely have missed the short-lived helper processes
that only remain as exited remnants in RAM. The Stuxnet image shows advanced persistence: kernel drivers and
fake credential-service processes that look like normal Windows names. Disk tools that only list running
processes by name would under-detect this. \textbf{Recommendation:} Add memory acquisition to Meridian's IR
playbook for ransomware and suspected credential/kernel compromise on legacy Windows assets, with hash
verification and dual process listing (active list plus pool scan) as mandatory steps.

\subsubsection*{What each image showed}
\textbf{WannaCry (\texttt{wcry.vmem}).}
 Interactive detonation under the user desktop process, staging under
\texttt{C:\textbackslash Intel\textbackslash\ldots}, encryptor process \texttt{@WanaDecryptor@}, and exited recovery-inhibition helpers.
SMB port~445 was listening on the system process (MS17-010 exposure). Capture time: 2017-05-12 21:26:32~UTC.

\textbf{Stuxnet (\texttt{stuxnet.vmem}).}
 One genuine Local Security Authority process (PID~680) versus two counterfeits
(PIDs~868 and~1928) created on 2011-06-03 under \texttt{services.exe}, with injected executable regions and
kernel drivers MRxCls/MRxNet registered seconds earlier. Capture time: 2011-06-03 04:31:36~UTC.

\subsubsection*{What memory found that disk-only triage would miss}
\begin{itemize}
  \item Exited WannaCry helpers (\texttt{taskse}/\texttt{taskdl}) with ExitTime still recoverable via pool scanning.
  \item Counterfeit \texttt{lsass.exe} instances that reuse a trusted name but fail parent/DLL/injection checks.
  \item Unlinked injected PE images (ldrmodules False/False/False) and RWX private VADs (malfind).
  \item Kernel driver load order relative to user-mode counterfeit process creation.
\end{itemize}

\subsubsection*{Consolidated IOC table (defanged)}
{\small
\begin{tabularx}{\textwidth}{@{}P{0.75in} Y P{0.95in} P{0.7in}@{}}
\toprule
Type & Indicator (defanged) & Source plugin & Conf. \\
\midrule
Process & tasksche.exe under \texttt{C:/Intel/.../tasksche.exe} & cmdline & High \\
Process & @WanaDecryptor@.exe & pslist/pstree & High \\
Process & taskse.exe / taskdl.exe (exited) & psscan/psxview & High \\
Network & TCP 445 listener (System PID~4) & sockets & High \\
Network & 192.168[.]56[.]101 & sockets & High \\
URL & hxxp://www[.]iuqerfsodp9ifjaposdfjhgosurijfaewrwergwea[.]com & strings/intel & Med \\
BTC & 12t9YDPgwueZ9NyMgw519p7AA8isjr6SMw & strings/intel & Med \\
Driver & mrxcls.sys / mrxnet.sys & driverscan & High \\
Hash & 2b2945f7\ldots ffe6d (malfind 0x80000) & malfind dump & High \\
Hash & 10f07b9f\ldots 3586 (VT Stuxnet/Duqu) & malfind + VT & Med \\
\bottomrule
\end{tabularx}}

\noindent\textbf{Investigated and assessed as benign (not dropped silently):}
{\small
\begin{tabularx}{\textwidth}{@{}P{1.35in} Y P{1.7in}@{}}
\toprule
Item & Why reviewed & Assessment \\
\midrule
Procmon (PID~660) & Unusual monitor on infection-day image & Lab tooling; not Stuxnet TTP \\
cmd / ipconfig (short-lived) & Appear on infection day & Ambiguous; not primary IOC \\
wuauclt / VMware tools & Present in trees & Expected guest/tools noise \\
Genuine lsass PID~680 & Name collision with counterfeits & Benign LSA (parent/dlllist/malfind) \\
\bottomrule
\end{tabularx}}

\subsubsection*{ATT\&CK table: Memory-observed vs Intel-sourced}
Technique IDs verified against attack.mitre.org software pages S0366 and S0603.

{\small
\begin{tabularx}{\textwidth}{@{}P{1.35in} P{0.55in} P{1.05in} Y@{}}
\toprule
Technique & ID & Evidence Source & Notes \\
\midrule
Data Encrypted for Impact & T1486 & Memory-observed & @WanaDecryptor@ / tasksche in wcry.vmem \\
Inhibit System Recovery & T1490 & Memory-observed & taskse/taskdl ExitTime (S0366) \\
Process Injection & T1055 & Memory-observed & malfind RWX MZ in fake lsass (S0603) \\
Rootkit & T1014 & Memory-observed & MRxCls/MRxNet loaded (S0603) \\
Exploitation for Client Execution & T1203 & Intel-sourced & MS17-010 not reconstructed from this RAM \\
Remote Services: SMB Admin Shares & T1021.002 & Intel-sourced & S0366 SMB worming; connscan empty here \\
\bottomrule
\end{tabularx}}

\subsubsection*{Three detection or prevention controls for Meridian}
\begin{enumerate}
\item \textbf{SMB isolation (signal: TCP/445 on System PID~4).}
 Detect/block at Meridian OT firewall and network ACL
between maintenance VLAN and plant HMIs. \textit{Limitation:} Host-only lab images prove exposure, not successful
lateral movement; empty connscan means this control must not wait for live C2 evidence.
\item \textbf{Dual process listing on IR (signal: pslist vs psscan ExitTime on taskse/taskdl).}
 Place in Meridian SOC
memory-triage playbook before rootkit escalation. \textit{Limitation:} Distinguishes exited helpers from DKOM hiding
only when ExitTime is checked; psxview False/True alone is insufficient.
\item \textbf{Module~02 WinPmem + Volatility~3 workflow.} Acquire with WinPmem, hash the raw dump, then run
\texttt{windows.info} $\rightarrow$ \texttt{windows.pslist} / \texttt{windows.psscan} $\rightarrow$ \texttt{windows.malfind} $\rightarrow$
\texttt{windows.dlllist} on candidate lsass PIDs and map hits to ATT\&CK. Place on Meridian IR workstation after isolation.
\textit{Limitation:} Symbol/profile coverage for XP-era images may force fallback to Volatility~2 profiles (as used here
with WinXPSP2x86).
\end{enumerate}

\subsubsection*{Business impact, legal, detection engineering, remediation}
For Meridian Energy, WannaCry-style ransomware in a host-only lab VLAN mirrors contractor laptops bridging
OT maintenance networks to corporate Wi-Fi. Encryption within minutes of execution (21:22 to 21:26~UTC) implies
backup and recovery playbooks must assume a blast radius under five minutes on unpatched SMB endpoints.
Stuxnet-style counterfeit lsass demonstrates credential subsystem compromise without obvious Process Explorer
anomalies; memory review is mandatory when EDR misses kernel-assisted injection.

Defanged IOCs in this report are suitable for external sharing; full paths under
\winpath{C:/Intel/...} remain internal. BTC addresses and onion endpoints are financial/extortion artifacts:
coordinate with law enforcement before blockchain tracing. Do not upload malfind dumps to public
multi-scanners from OT environments (B9 OPSEC).

Priority detection candidates: (1)~child process explorer $\rightarrow$ unsigned tasksche in a non-standard directory;
(2)~multiple lsass.exe with the same ImageFileName but divergent dlllist counts;
(3)~kernel driver load MRxCls+MRxNet paired with new lsass instances;
(4)~RWX VADs with MZ headers inside lsass (malfind signature).

\textbf{Remediation sequence:} Isolate $\rightarrow$ image memory $\rightarrow$ hash verify $\rightarrow$ Volatility profile $\rightarrow$
pslist/psscan delta $\rightarrow$ cmdline/dlllist on suspicious PIDs $\rightarrow$ malfind dump $\rightarrow$ hash-only VT $\rightarrow$
reimage from known-good baseline. For ICS, validate PLC project backups out-of-band; RAM artifacts do not
prove PLC payload delivery (C1 claim~6).

\textbf{Open questions:} (1)~Whether wcry.vmem captured before or after full disk encryption completion;
(2)~provenance of missing PPID~1608; (3)~whether Procmon in stuxnet.vmem alters timing of counterfeit lsass;
(4)~services.exe malfind region 0x13f0000 needs deeper PE carve for attribution to Stuxnet versus tooling.

\subsection*{D2. Knowledge check: True/False with corrections}
{\footnotesize
\begin{tabularx}{\textwidth}{@{}c l >{\RaggedRight\arraybackslash}Y@{}}
\toprule
\# & Verdict & Correction / rationale \\
\midrule
1 & Correct & pslist walks the linked EPROCESS list; psscan searches pool tags (different visibility). \\
2 & Incorrect & psxview False/True for taskse/taskdl is ExitTime exited helpers, not DKOM rootkit hiding. \\
3 & Incorrect & connscan/sockets are point-in-time; empty WannaCry PID rows do not prove no SMB spread. \\
4 & Incorrect & imageinfo +0530 is victim host TZ configuration, not attacker geolocation (India). \\
5 & Incorrect & PID~680 owning UDP 500/4500 is corroborative only; genuineness needs parent/dlllist/malfind. \\
6 & Incorrect & \texttt{dlllist | wc -l} counts banners/headers; scoped 0x-row counts are 57 (680) / 8 (868) / 28 (1928). \\
7 & Incorrect & DEC EBP; POP EDX; NOP is MZ/PE header mis-disassembly, not Stuxnet injection shellcode. \\
8 & Incorrect & malfind -D yields one file per VAD; six unique SHA-256s; VT 61/71 Stuxnet vs 48/70 Duqu. \\
9 & Correct (nuance) & WinXPSP2x86 is the validated Volatility~2 profile even when Image Type reports SP3. \\
10 & Correct & Create/Exit fields print UTC+0000; imageinfo separately shows host-local offset. \\
\bottomrule
\end{tabularx}}

\section*{Analytical Standards}
Evidence tiers: \textbf{Tier~1} direct Volatility plugin output; \textbf{Tier~2} multi-plugin agreement
(pslist/psscan/ExitTime); \textbf{Tier~3} dossier/ATT\&CK not independently observed in RAM.
Profile discipline: imageinfo $\rightarrow$ instantiate WinXPSP2x86 $\rightarrow$ validate with pslist.
WannaCry focus: exited versus hidden (ExitTime). Stuxnet focus: counterfeit lsass via parent, dlllist, malfind, sockets.
Hygiene: defanged IOCs, hash-only VT, UTC timelines, attribution limits.


\section*{Evidence Appendix}
\subsection*{Evidence log (summary)}
Intake and final hashes match (see intake table). Full START/DONE command log with timestamps is retained under
\texttt{work/output/command\_log.txt}. SIFT confirmation outputs are under \texttt{work/output/sift/}.

\begin{figure}[H]
\centering
\includegraphics[width=0.98\textwidth]{figures/fig_SIFT_login.png}
\caption{SIFT Workstation login (sansforensics) after OVA import into VirtualBox.}
\end{figure}\begin{figure}[H]
\centering
\includegraphics[width=0.98\textwidth]{figures/fig_SIFT_wcry_pslist.png}
\caption{SIFT GUI terminal: vol.py pslist on wcry.vmem showing tasksche.exe (PID 1940) and @WanaDecryptor@ (PID 740).}
\end{figure}\begin{figure}[H]
\centering
\includegraphics[width=0.98\textwidth]{figures/fig_SIFT_stux_lsass.png}
\caption{SIFT vol.py pslist on stuxnet.vmem (lsass.exe rows) confirming the triple-lsass process set.}
\end{figure}
\noindent\textbf{Command log (excerpt)}
\begin{lstlisting}
MIDTERM VOL2 ANALYSIS 2026-10-05T16:28:31.4407818-04:00
[2026-10-05T16:28:31.4653891-04:00] START C:\temp\midterm\tools\vol26.exe -f C:\temp\midterm\dumps\wcry.vmem imageinfo
[2026-10-05T16:28:31.4653891-04:00] DONE lines=12
[2026-10-05T16:28:45.7628248-04:00] START C:\temp\midterm\tools\vol26.exe -f C:\temp\midterm\dumps\wcry.vmem kdbgscan
[2026-10-05T16:28:45.7628248-04:00] DONE lines=32
[2026-10-05T16:28:57.9841232-04:00] START C:\temp\midterm\tools\vol26.exe -f C:\temp\midterm\dumps\stuxnet.vmem imageinfo
[2026-10-05T16:28:57.9841232-04:00] DONE lines=12
[2026-10-05T16:29:22.1335656-04:00] START C:\temp\midterm\tools\vol26.exe -f C:\temp\midterm\dumps\stuxnet.vmem kdbgscan
[2026-10-05T16:29:22.1335656-04:00] DONE lines=32
[2026-10-05T16:29:40.4588895-04:00] START C:\temp\midterm\tools\vol26.exe -f C:\temp\midterm\dumps\wcry.vmem --profile WinXPSP2x86 pslist
[2026-10-05T16:29:40.4588895-04:00] DONE lines=22
[2026-10-05T16:29:42.3355904-04:00] START C:\temp\midterm\tools\vol26.exe -f C:\temp\midterm\dumps\wcry.vmem --profile WinXPSP2x86 pstree
[2026-10-05T16:29:42.3355904-04:00] DONE lines=22
[2026-10-05T16:29:44.4015077-04:00] START C:\temp\midterm\tools\vol26.exe -f C:\temp\midterm\dumps\wcry.vmem --profile WinXPSP2x86 psscan
[2026-10-05T16:29:44.4015077-04:00] DONE lines=27
[2026-10-05T16:29:46.4772117-04:00] START C:\temp\midterm\tools\vol26.exe -f C:\temp\midterm\dumps\wcry.vmem --profile WinXPSP2x86 psxview
[2026-10-05T16:29:46.4772117-04:00] DONE lines=27
[2026-10-05T16:29:49.7528941-04:00] START C:\temp\midterm\tools\vol26.exe -f C:\temp\midterm\dumps\wcry.vmem --profile WinXPSP2x86 connections
[2026-10-05T16:29:49.7528941-04:00] DONE lines=2
[2026-10-05T16:29:51.4218671-04:00] START C:\temp\midterm\tools\vol26.exe -f C:\temp\midterm\dumps\wcry.vmem --profile WinXPSP2x86 connscan
[2026-10-05T16:29:51.4218671-04:00] DONE lines=2
[2026-10-05T16:29:53.9006297-04:00] START C:\temp\midterm\tools\vol26.exe -f C:\temp\midterm\dumps\wcry.vmem --profile WinXPSP2x86 sockets
[2026-10-05T16:29:53.9006297-04:00] DONE lines=19
[2026-10-05T16:29:55.5530276-04:00] START C:\temp\midterm\tools\vol26.exe -f C:\temp\midterm\dumps\wcry.vmem --profile WinXPSP2x86 sockscan
[2026-10-05T16:29:55.5530276-04:00] DONE lines=19
[2026-10-05T16:29:57.7298516-04:00] START C:\temp\midterm\tools\vol26.exe -f C:\temp\midterm\dumps\wcry.vmem --profile WinXPSP2x86 cmdline
[2026-10-05T16:29:57.7298516-04:00] DONE lines=59
[2026-10-05T16:29:59.5509123-04:00] START C:\temp\midterm\tools\vol26.exe -f C:\temp\midterm\dumps\wcry.vmem --profile WinXPSP2x86 consoles
[2026-10-05T16:29:59.5509123-04:00] DONE lines=30
[2026-10-05T16:30:01.9208025-04:00] START C:\temp\midterm\tools\vol26.exe -f C:\temp\midterm\dumps\wcry.vmem --profile WinXPSP2x86 envars
[2026-10-05T16:30:01.9208025-04:00] DONE lines=414
[2026-10-05T16:30:03.9072806-04:00] START C:\temp\midterm\tools\vol26.exe -f C:\temp\midterm\dumps\wcry.vmem --profile WinXPSP2x86 dlllist
[2026-10-05T16:30:03.9072806-04:00] DONE lines=1020
[2026-10-05T16:30:05.8255541-04:00] START C:\temp\midterm\tools\vol26.exe -f C:\temp\midterm\dumps\stuxnet.vmem --profile WinXPSP2x86 pslist
[2026-10-05T16:30:05.8255541-04:00] DONE lines=33
... (82 lines total; truncated)
\end{lstlisting}

\noindent\textbf{WannaCry psscan (exited helpers)}
\begin{lstlisting}
Offset(P)          Name                PID   PPID PDB        Time created                   Time exited
------------------ ---------------- ------ ------ ---------- ------------------------------ ------------------------------
0x0000000001f4daf0 taskdl.exe          860   1940 0x199f6000 2017-05-12 21:26:23 UTC+0000   2017-05-12 21:26:23 UTC+0000
0x0000000001f53d18 taskse.exe          536   1940 0x1986c000 2017-05-12 21:26:22 UTC+0000   2017-05-12 21:26:23 UTC+0000
0x0000000001f69b50 @WanaDecryptor@     424   1940 0x18fa2000 2017-05-12 21:25:52 UTC+0000   2017-05-12 21:25:53 UTC+0000
0x0000000001f747c0 wuauclt.exe        1768   1024 0x11629000 2017-05-12 21:22:52 UTC+0000
0x0000000001f8ba58 @WanaDecryptor@     576   1940 0x19671000 2017-05-12 21:26:22 UTC+0000   2017-05-12 21:26:23 UTC+0000
0x0000000001fb95d8 svchost.exe         260    664 0x0ce48000 2017-05-12 21:22:18 UTC+0000
0x0000000001fde308 @WanaDecryptor@     740   1940 0x0de3a000 2017-05-12 21:22:22 UTC+0000
0x0000000001fea8a0 wscntfy.exe        1168   1024 0x12217000 2017-05-12 21:22:56 UTC+0000
0x0000000001ffa710                       0      0 0x17d3f000
0x0000000002010020 alg.exe             544    664 0x1238d000 2017-05-12 21:22:55 UTC+0000
0x000000000203b7a8 svchost.exe        1084    664 0x0838c000 2017-05-12 21:22:03 UTC+0000
0x0000000002161da0 csrss.exe           596    348 0x07752000 2017-05-12 21:22:00 UTC+0000
0x0000000002169020 smss.exe            348      4 0x0683e000 2017-05-12 21:21:55 UTC+0000
0x000000000216e020 winlogon.exe        620    348 0x07957000 2017-05-12 21:22:01 UTC+0000
0x0000000002191658 lsass.exe           676    620 0x07bb7000 2017-05-12 21:22:01 UTC+0000
... (27 lines total; truncated)
\end{lstlisting}

\noindent\textbf{Stuxnet dlllist PID 868}
\begin{lstlisting}
************************************************************************
lsass.exe pid:    868
Command line : "C:\WINDOWS\\system32\\lsass.exe"
Service Pack 3

Base             Size  LoadCount Path
---------- ---------- ---------- ----
0x01000000     0x6000     0xffff C:\WINDOWS\system32\lsass.exe
0x7c900000    0xaf000     0xffff C:\WINDOWS\system32\ntdll.dll
0x7c800000    0xf6000     0xffff C:\WINDOWS\system32\kernel32.dll
0x77dd0000    0x9b000     0xffff C:\WINDOWS\system32\ADVAPI32.dll
0x77e70000    0x92000     0xffff C:\WINDOWS\system32\RPCRT4.dll
0x77fe0000    0x11000     0xffff C:\WINDOWS\system32\Secur32.dll
0x7e410000    0x91000     0xffff C:\WINDOWS\system32\USER32.dll
0x77f10000    0x49000     0xffff C:\WINDOWS\system32\GDI32.dll
\end{lstlisting}

\noindent\textbf{Stuxnet malfind PID 1928}
\begin{lstlisting}
Process: lsass.exe Pid: 1928 Address: 0x80000
Vad Tag: Vad  Protection: PAGE_EXECUTE_READWRITE
Flags: Protection: 6

0x00080000  4d 5a 90 00 03 00 00 00 04 00 00 00 ff ff 00 00   MZ..............
0x00080010  b8 00 00 00 00 00 00 00 40 00 00 00 00 00 00 00   ........@.......
0x00080020  00 00 00 00 00 00 00 00 00 00 00 00 00 00 00 00   ................
0x00080030  00 00 00 00 00 00 00 00 00 00 00 00 08 01 00 00   ................

0x00080000 4d               DEC EBP
0x00080001 5a               POP EDX
0x00080002 90               NOP
0x00080003 0003             ADD [EBX], AL
0x00080005 0000             ADD [EAX], AL
0x00080007 000400           ADD [EAX+EAX], AL
0x0008000a 0000             ADD [EAX], AL
0x0008000c ff               DB 0xff
... (197 lines total; truncated)
\end{lstlisting}

\noindent\textbf{Stuxnet modules (mrx*)}
\begin{lstlisting}
Offset(V)  Name                 Base             Size File
---------- -------------------- ---------- ---------- ----
0x823fc3a0 ntoskrnl.exe         0x804d7000   0x1f8580 \WINDOWS\system32\ntkrnlpa.exe
0x823fc338 hal.dll              0x806d0000    0x20300 \WINDOWS\system32\hal.dll
0x823fc2d0 kdcom.dll            0xf8b9a000     0x2000 \WINDOWS\system32\KDCOM.DLL
0x823fc260 BOOTVID.dll          0xf8aaa000     0x3000 \WINDOWS\system32\BOOTVID.dll
0x823fc1f8 ACPI.sys             0xf856b000    0x2e000 ACPI.sys
0x823fc188 WMILIB.SYS           0xf8b9c000     0x2000 \WINDOWS\system32\DRIVERS\WMILIB.SYS
0x823fc120 pci.sys              0xf855a000    0x11000 pci.sys
0x823fc0b0 isapnp.sys           0xf869a000     0xa000 isapnp.sys
0x823fc040 compbatt.sys         0xf8aae000     0x3000 compbatt.sys
0x823ed008 BATTC.SYS            0xf8ab2000     0x4000 \WINDOWS\system32\DRIVERS\BATTC.SYS
0x823edf98 intelide.sys         0xf8b9e000     0x2000 intelide.sys
0x823edf28 PCIIDEX.SYS          0xf891a000     0x7000 \WINDOWS\system32\DRIVERS\PCIIDEX.SYS
0x823edeb8 MountMgr.sys         0xf86aa000     0xb000 MountMgr.sys
... (124 lines total; truncated)
\end{lstlisting}

\section*{Method Map (six-step workflow, course slide 17)}
One-page map placing Part~A/~B findings under the six investigation stages.

{\footnotesize
\begin{tabularx}{\textwidth}{@{}c l >{\RaggedRight\arraybackslash}Y@{}}
\toprule
Step & Stage & Findings placed here \\
\midrule
1 & Acquire \& hash & Intake SHA-256 for wcry.vmem, stuxnet.vmem, zip, dossier; chain of custody \\
2 & Profile & A1/B1 imageinfo; WinXPSP2x86; PAE difference (wcry No / stuxnet Yes); KDBG/DTB \\
3 & Process review & A2 pstree lineage; A3/A4 psxview ExitTime; B2/B3 genuine vs counterfeit lsass; B5/B6 dlllist \\
4 & Network & A5 sockets/connscan (445 / empty WannaCry rows); B4 UDP 500/4500 on PID 680 only \\
5 & Deep dive & A7 cmdline paths; B7 malfind; B8 dumps/hashes; B9 VT hash-only; B10 ldrmodules/drivers \\
6 & Correlate \& report & A4/A8 and D1 ATT\&CK Memory-observed vs Intel-sourced; C1 to C3 dossier trial; D1 memo \\
\bottomrule
\end{tabularx}}

\section*{AI Disclosure}
Generative AI assisted with report drafting, LaTeX structuring, VirtualBox/SIFT bring-up, and Volatility
command-syntax sanity checks. All forensic values (hashes, PIDs, timestamps, malfind regions, dll counts,
socket bindings, and Volatility excerpts) originate from the analyst's own Volatility~2.6 runs
(\texttt{vol26.exe} on the host and \texttt{vol.py} inside SIFT) on local copies of \texttt{wcry.vmem} and
\texttt{stuxnet.vmem} captured October~5, 2026. AI did not replace hands-on plugin execution or hash verification.

\section*{References (APA 7)}
\begin{enumerate}[leftmargin=1.6em,itemsep=0.35em]
\item Ligh, M.~H., Case, A., Levy, J., \& Walters, M. (2014). \textit{The art of memory forensics: Detecting malware and threats in Windows, Linux, and Mac memory}. Wiley.
\item MITRE ATT\&CK. (n.d.). \textit{WannaCry} (Software S0366). \url{https://attack.mitre.org/software/S0366/}
\item MITRE ATT\&CK. (n.d.). \textit{Stuxnet} (Software S0603). \url{https://attack.mitre.org/software/S0603/}
\item Microsoft. (2017). Microsoft security bulletin MS17-010. \url{https://learn.microsoft.com/en-us/security-updates/SecurityBulletins/2017/ms17-010}
\item Symantec Security Response. (2011). \textit{W32.Stuxnet Dossier} (Version 1.3). Symantec Corp. (Sections cited: Executive Summary; Injection Technique; Installation; Load Point; Bypassing Behavior Blocking When Loading DLLs; Windows Rootkit Functionality.)
\item The Volatility Foundation. (n.d.). Volatility 2.6 documentation.
\item VirusTotal. (2026). Hash intelligence lookups performed 2026-10-05 (query-only, no sample upload).
\end{enumerate}

\end{document}
